\section{Quadrature}

Numerical quadrature aims to compute an approximation of a definite integral:
$$I(f) := \int_a^b f(x) \, dx$$
using a finite weighted sum of function evaluations:
$$Q(f) := \sum_{j=0}^n w_j f(c_j)$$
where $c_j \in [a, b]$ are the \textbf{quadrature nodes} and $w_j \in \mathbb{R}$ are the \textbf{quadrature weights}.

\subsection{Polynomial Interpolation}

\begin{definition}
    \textbf{Interpolatory Quadrature via Lagrange Basis:} To approximate the integral of an unknown or complicated function $f(x)$, we interpolate $f$ at $n+1$ distinct nodes $c_0, c_1, \dots, c_n \in [a,b]$ with a polynomial $p_n \in \mathcal{P}_n$, and integrate $p_n$ analytically:
    $$I(f) \approx \int_a^b p_n(x) \, dx$$
    Expressing the interpolant $p_n(x)$ in the \textbf{Lagrange basis} $L_j(x) = \prod_{k \neq j} \frac{x - c_k}{c_j - c_k}$:
    $$p_n(x) = \sum_{j=0}^n f(c_j) L_j(x)$$
    Integrating term-by-term yields:
    $$\int_a^b p_n(x) \, dx = \sum_{j=0}^n \left( \int_a^b L_j(x) \, dx \right) f(c_j)$$
    Thus, the quadrature weights are uniquely determined by integrating the Lagrange basis polynomials:
    $$w_j = \int_a^b L_j(x) \, dx \quad \text{for } j = 0, 1, \dots, n$$
\end{definition}

\begin{remark}
    An interpolatory quadrature rule with $n+1$ nodes has degree of exactness at least $n$: it integrates every polynomial $p \in \mathcal{P}_n$ exactly, because polynomial interpolation reproduces polynomials of degree $\le n$ without error ($p_n \equiv p$).
\end{remark}

\begin{important}
    \textbf{Runge's Phenomenon in Quadrature:} Just as high-degree polynomial interpolation on \textbf{equidistant nodes} suffers from Runge's phenomenon (wild oscillations near the boundary), using equidistant nodes for high-order quadrature leads to severe instability.
    For $n \ge 8$, the weights $w_j = \int_a^b L_j(x) dx$ begin to take on alternating negative signs, and:
    $$\sum_{j=0}^n |w_j| \xrightarrow{n \to \infty} \infty \quad \text{exponentially}$$
    Since the condition number of quadrature is $\sum |w_j|$, large opposing weights lead to \textbf{catastrophic numerical cancellation}. Therefore, high-degree polynomial interpolation on equidistant nodes should never be used for quadrature.
\end{important}

\subsection{Composite Quadrature Rules}

\begin{definition}
    \textbf{Domain Partitioning ($h$-Refinement):} Instead of increasing the polynomial degree $n$ over the whole interval, we keep the degree low and divide $[a,b]$ into $N$ smaller subintervals of uniform width $h = \frac{b-a}{N}$:
    $$a = x_0 < x_1 < x_2 < \dots < x_N = b, \quad \text{where } x_i = a + i h$$
    The total integral is the sum of integrals over each subinterval $[x_{i-1}, x_i]$:
    $$\int_a^b f(x) \, dx = \sum_{i=1}^N \int_{x_{i-1}}^{x_i} f(x) \, dx$$
\end{definition}

\subsubsection{The Midpoint Rule}

\begin{definition}
    \textbf{Midpoint Rule:} Approximates each subinterval's integral by the area of a rectangle whose height is the function value at the subinterval midpoint $x_{i-1} + \frac{1}{2}h$:
    $$Q_N^{MP}(f) := h \sum_{i=1}^N f\left(x_{i-1} + \frac{1}{2}h\right)$$
\end{definition}

\begin{theorem}
    \textbf{Midpoint Error:} For $f \in C^2[a,b]$, the composite midpoint error satisfies:
    $$|I(f) - Q_N^{MP}(f)| = \frac{b-a}{24} h^2 |f''(\xi)| = \mathcal{O}(h^2) \quad \text{for some } \xi \in [a,b]$$
\end{theorem}

\begin{remark}
    \textbf{Intuition for the Midpoint "Bonus Degree":} Although the midpoint rule approximates $f$ with a degree-0 constant, it integrates linear functions $f(x) = c_1 x + c_0$ exactly! By symmetry around the midpoint $m$, the odd linear error term $(x-m)$ integrates to zero: the triangular error above the midpoint rectangle exactly cancels the triangular error below it.
\end{remark}

\subsubsection{The Trapezoidal Rule}

\begin{definition}
    \textbf{Trapezoidal Rule:} Approximates each subinterval's integral as the trapezoid spanned by its endpoints:
    \begin{align*}
        Q_N^{Trap}(f) &:= \sum_{i=1}^N \frac{h}{2} \Big(f(x_{i-1}) + f(x_i)\Big) \\
        &= h \left[ \frac{1}{2}f(x_0) + f(x_1) + f(x_2) + \dots + f(x_{N-1}) + \frac{1}{2}f(x_N) \right]
    \end{align*}
    Every interior node is shared between adjacent intervals and receives weight $h$, while boundary endpoints receive weight $\frac{h}{2}$.
\end{definition}

\begin{theorem}
    \textbf{Trapezoidal Error:} For $f \in C^2[a,b]$, the composite trapezoidal error satisfies:
    $$|I(f) - Q_N^{Trap}(f)| = \frac{b-a}{12} h^2 |f''(\eta)| = \mathcal{O}(h^2) \quad \text{for some } \eta \in [a,b]$$
\end{theorem}

\subsubsection{Simpson's Rule}

\begin{definition}
    \textbf{Simpson's Rule:} Interpolates $f$ with a quadratic parabola over each subinterval using the two endpoints and the midpoint, weighted as $\frac{h}{6}(1, 4, 1)$:
    $$Q_N^{Simp}(f) := \frac{h}{6} \sum_{i=1}^N \left[ f(x_{i-1}) + 4 f\left(\frac{x_{i-1}+x_i}{2}\right) + f(x_i) \right]$$
\end{definition}

\begin{theorem}
    \textbf{Simpson's Error:} Due to point symmetry around the midpoint, the cubic error term $(x-m)^3$ cancels out, making Simpson's rule exact for all cubic polynomials $\mathcal{P}_3$. For $f \in C^4[a,b]$:
    $$|I(f) - Q_N^{Simp}(f)| = \frac{b-a}{2880} h^4 |f^{(4)}(\xi)| = \mathcal{O}(h^4)$$
    Halving the step size $h \to \frac{h}{2}$ reduces the error by a factor of $2^4 = 16$.
\end{theorem}

\subsubsection{Smoothness}

\begin{remark}
    \textbf{Non-Smooth Functions Degrade Convergence:} The theoretical error bounds $\mathcal{O}(h^2)$ and $\mathcal{O}(h^4)$ strictly require bounded continuous derivatives ($f \in C^2$ or $f \in C^4$).
    \begin{itemize}
        \item If $f$ has a \textbf{kink} (continuous, but first derivative jumps, e.g. $f(x) = |x - c|$), $f''$ does not exist as a bounded function. The convergence rate of both Trapezoidal and Simpson drops to $\mathcal{O}(h^1)$.
        \item If $f$ has a \textbf{singularity} or infinite slope (e.g. $f(x) = \sqrt{x}$ near 0), the convergence rate degrades further to $\mathcal{O}(h^{0.5})$.
    \end{itemize}
    High-order rules offer no advantage when the integrand lacks smoothness!
\end{remark}


\subsubsection{Romberg Extrapolation}

\begin{remark}
    textbf{Intuition:} By knowing the mathematical structure of the error, we can combine two relatively poor approximations to eliminate the dominant error term, creating a much higher-order method for free!.
\end{remark}

\begin{definition}
    \textbf{Romberg Extrapolation for Quadrature:} Applying Richardson extrapolation to the composite trapezoidal rule is called \textbf{Romberg Integration}.
    By the Euler-Maclaurin expansion, the trapezoidal error consists \textbf{exclusively of even powers of $h$}:
    $$T(h) = I(f) + c_2 h^2 + c_4 h^4 + c_6 h^6 + c_8 h^8 + \dots$$
    Note that intuitively this comes from the trapezoid being symetric hence $T(h) = T(-h)$, if there were odd terms this would not hold because of the changing sign.
    And because there are no odd powers ($h^1, h^3, h^5$), the leading error term has $p = 2$.
    $$T_1(h/2) = \frac{4 T(h/2) - T(h)}{4 - 1} = \frac{4 T(h/2) - T(h)}{3}$$
    This cancels the $c_2 h^2$ error term, yielding an $\mathcal{O}(h^4)$ method.
    \emph{Remarkable fact:} $T_1(h/2)$ is algebraically identical to \textbf{Simpson's rule}!
\end{definition}

\subsubsection{High-Order Rules}

\begin{remark}
    High-order \textbf{Newton-Cotes} rules fit a single high-degree polynomial over equidistant nodes on the entire interval. As established by Runge's phenomenon, this leads to negative weights for $n \ge 8$ and exponential divergence of condition numbers ($\sum |w_j| \to \infty$). Hence, high-degree equidistant rules are unstable and should be avoided.
\end{remark}

\begin{definition}
    \textbf{Gauss Quadrature (Gauss-Legendre):} Instead of fixing equidistant nodes, Gauss quadrature treats \textbf{both} the $n+1$ nodes $c_j$ and the $n+1$ weights $w_j$ as free variables ($2n+2$ parameters).
    \begin{itemize}
        \item \textbf{Maximal Exactness:} Achieves the maximum theoretically possible degree of exactness:
        $$d = 2n + 1 \quad \iff \quad \text{Order } q = 2n + 2$$
        \item \textbf{Orthogonality:} The nodes $c_0, \dots, c_n$ must be the $n+1$ roots of the $(n+1)$-th \textbf{Legendre polynomial} $P_{n+1}(x)$ on $[-1, 1]$.
        \item \textbf{Positivity \& Stability:} All weights are strictly positive ($w_j > 0$), guaranteeing unconditional stability ($\sum |w_j| = b-a$).
        \item \textbf{Drawback:} Gauss nodes are \textbf{non-nested}: the nodes for $n$ and $n+1$ do not overlap, so changing order requires throwing away all previous function evaluations.
    \end{itemize}
\end{definition}

\begin{definition}
    \textbf{Clenshaw-Curtis Quadrature:} Clenshaw-Curtis quadrature interpolates $f$ at the \textbf{Chebyshev-Gauss-Lobatto nodes} (extrema of Chebyshev polynomials, which include the endpoints):
    $$x_j = \cos\left(\frac{j\pi}{n}\right) \quad \text{for } j = 0, 1, \dots, n$$
    and integrates the Chebyshev polynomial interpolant analytically.
\end{definition}

\begin{remark}
    \textbf{Key Advantages of Clenshaw-Curtis:}
    \begin{enumerate}
        \item \textbf{Suppresses Runge's Phenomenon:} Because Chebyshev nodes cluster near the boundaries with density $\propto \frac{1}{\sqrt{1-x^2}}$, the polynomial interpolation error is well-behaved and weights remain strictly positive for all practical $n$.
        \item \textbf{Nested Nodes:} When doubling the number of intervals $n \to 2n$, all previous nodes are retained:
        $$x_j^{(n)} = \cos\left(\frac{j\pi}{n}\right) = \cos\left(\frac{2j\pi}{2n}\right) = x_{2j}^{(2n)}$$
        This allows adaptive doubling without wasting any prior function evaluations!
        \item \textbf{Fast Computation via FFT:} The quadrature weights and evaluations can be computed in only $\mathcal{O}(n \log n)$ operations using the Fast Cosine Transform (FFT).
        \item \textbf{Near-Gaussian Performance:} In practice, for smooth functions, Clenshaw-Curtis converges almost as quickly as Gauss quadrature while offering nested convenience.
    \end{enumerate}
\end{remark}

\subsection{Adaptive Quadrature}

\begin{remark}
    \textbf{Motivation for Adaptivity:} When an integrand $f$ has a localized sharp spike, steep wave, or near-singularity, a uniform step size $h$ is severely wasteful. To resolve the spike, $h$ must be tiny everywhere, wasting computation in smooth regions. \textbf{Adaptive quadrature} refines $h$ only where needed.
\end{remark}

\begin{definition}
    \textbf{Local Error Estimation via Step Halving:} On any subinterval $[a, b]$ with midpoint $m = \frac{a+b}{2}$:
    \begin{itemize}
        \item Compute a coarse estimate $I_{\text{coarse}} = Q(a, b)$ (e.g., single Simpson step of width $h = b-a$).
        \item Compute a fine estimate $I_{\text{fine}} = Q(a, m) + Q(m, b)$ (two Simpson steps of width $h/2$).
        \item Estimate the local truncation error:
        $$\text{Est}(a, b) := \frac{1}{15} |I_{\text{fine}} - I_{\text{coarse}}| \approx |I(f) - I_{\text{fine}}|$$
        (where the factor $\frac{1}{15} = \frac{1}{2^4 - 1}$ comes directly from Richardson extrapolation for Simpson's rule!).
    \end{itemize}
\end{definition}

\begin{definition}
    \textbf{Recursive Adaptive Algorithm:} Given a user-specified global tolerance $\text{tol} > 0$:
    \begin{enumerate}
        \item If $\text{Est}(a, b) \le \text{tol} \cdot \frac{b-a}{B-A}$ (where $[B, A]$ is the global domain):
        The local error is acceptable; return $I_{\text{fine}}$.
        \item If $\text{Est}(a, b) > \text{tol} \cdot \frac{b-a}{B-A}$:
        Subdivide $[a, b]$ into $[a, m]$ and $[m, b]$, and recursively apply the algorithm to each half with local tolerance $\text{tol}/2$.
    \end{enumerate}
\end{definition}

\begin{remark}
    \textbf{Intuition:} Adaptive quadrature acts like a dynamic magnifying glass: it recursively zooms in with dense evaluations around sharp spikes and steep gradients, while taking large, efficient steps across flat, smooth regions.
\end{remark}
